\begin{textsample}{NEG Dim 1 – 2011\_c4 – Score -4.00 – 2011\_c4/enade\_2011\_tec\_sane\_amb\_4989e09b77.txt}  \label{ex:f1_neg_114}
CAPÍTULO 4 PERCEPÇÃO DA PROVA As análises feitas neste capítulo tratam das percepções dos concluintes da Área de Tecnologia em Saneamento Ambiental sobre a prova aplicada no ENADE/2011.
Estas percepções foram mensuradas por meio de nove questões que avaliaram desde o grau de dificuldade da prova até o tempo gasto para concluí -la.
As percepções sobre a prova foram relacionadas com o desempenho dos estudantes e com a Grande Região de funcionamento do curso.
O questionário de percepção da prova encontra -se ao final do Anexo IV, que traz a reprodução da prova.
O desempenho dos estudantes foi classificado em quatro quartos.
Para tanto, esse desempenho foi ordenado de forma ascendente.
O percentil 25, P25, também conhecido como primeiro quartil, é a nota de desempenho que deixa um quarto ( 25 \% ) dos valores observados abaixo e três quartos acima.
A Figura 1 apresenta uma ilustração deste conceito.
O quarto inferior de desempenho é composto pelas notas abaixo do primeiro quartil.
Já o percentil 75, P75, também conhecido como terceiro quartil é o valor para o qual há três quartos ( 75 \% ) dos dados abaixo e um quarto acima dele.
O quarto superior de desempenho é composto pelas notas iguais ou acima do terceiro quartil.
O percentil 50, P50, também conhecido como mediana, é o valor que divide as notas em dois conjuntos de igual tamanho.
O segundo quarto inclui valores entre o primeiro quartil ( P25 ) e a mediana.
O terceiro quarto contém os valores entre a mediana ( P50 ) e o terceiro quartil ( P75 ).
Vale ressaltar que percentis, quartis e medianas são pontos que não obrigatoriamente pertencem ao conjunto original de dados, ao passo que os quartos são subconjuntos dos dados originais.
A seguir, serão apresentados gráficos com resultados selecionados, relativos às nove questões avaliadas por grupos de estudantes.
Os gráficos apresentam nas barras o percentual de alunos que assinalaram uma das opções ou a soma das porcentagens daqueles que assinalaram duas ( ou três ) delas.
Por exemplo, para as questões 1 e 2, os gráficos apresentam a porcentagem total de participantes que assinalaram as opções ( D ) difícil e ( E ) muito difícil.
Em \textbf{cada} barra foram assinalados também os extremos do intervalo de confiança de 95 \% como linhas verticais unidas por uma linha horizontal na forma da letra H maiúscula, semelhantemente aos gráficos do Capítulo 3.
As Tabelas no Anexo II apresentam os valores absolutos e a distribuição percentual das alternativas válidas das nove questões, segundo o mesmo recorte de desempenho dos alunos e Grande Região de funcionamento do curso. 4.1 GRAU DE DIFICULDADE DA PROVA 4.1.1 Componente de Formação Geral Ao avaliarem “ Qual o grau de dificuldade desta prova na parte de Formação Geral? ” ( Questão 1 ), 12,5 \% do grupo de inscritos e presentes optaram pelas alternativas difícil ou muito difícil.
Entretanto, para mais da metade dos estudantes ( 65,7 \% ) o Componente de Formação Geral da prova foi considerado com grau de dificuldade médio ( Gráfico 4.1, Gráfico 4.2 e, no Anexo II, a Tabela II.1 ).
O percentual de estudantes que consideraram a prova como difícil ou muito difícil foi maior na região Nordeste, onde a proporção foi de 17,0 \%, enquanto nas regiões nas regiões Sul e a Centro-Oeste nenhum estudante considerou a prova difícil ou muito difícil.
No Gráfico 4.1 é possível observar que as diferenças entre as regiões em que houve alunos que responderam que a parte da prova relativa ao Componente de Formação Geral foi difícil ou muito difícil não são estatisticamente significativas.
Nas Grandes Regiões, a proporção de presentes à prova que consideraram o Componente de Formação Geral como sendo de grau de dificuldade médio \textbf{esteve} entre 53,1 \% na região Centro-Oeste e 100,0 \% na região Sul.
O \textbf{menor} percentual de alunos que consideraram a prova difícil ou muito difícil ocorreu no quarto de desempenho superior ( 4º quarto ), onde foi de 6,9 \%.
A maior proporção foi 17,9 \%, dos estudantes do 2º quarto.
Com diferenças que não chegaram a ser significas estatisticamente, nota -se uma tendência decrescente do 2º para o 4º quartos.
Já entre o 1º e o 2º quartos a proporção de alunos que consideraram a prova difícil ou muito difícil aumenta.
Em todos os quartos a alternativa modal para esta pergunta foi médio, variando de 59,8 \% ( 4º quarto ) a 72,0 \% ( 1º quarto ). 4.1.2 Componente de Conhecimento Específico Ao responderem à Questão 2 – “ Qual o grau de dificuldade desta prova na parte de Componente Específico? ” – 29,7 \% dos estudantes da Área de tecnologia em Saneamento Ambiental classificaram -na como difícil ou muito difícil.
Além disso, o Componente de Conhecimento Específico da prova foi considerado com grau de dificuldade médio por 63,1 \% dos alunos ( Gráfico 4.3, Gráfico 4.4, e, no Anexo II, a Tabela II.2 ).
A maioria dos estudantes, de todas as Grandes Regiões, respondeu que o grau de dificuldade do Componente de Conhecimento Específico da prova era médio, com proporções variando de 54,5 \% a 67,7 \%, para regiões Norte e Sudeste, respectivamente.
A diferença entre a maior e a menor proporção de alunos que a avaliaram o grau de dificuldade como difícil ou muito difícil não é estatisticamente significativa : 42,4 \% na Norte e 21,9 \% na região Sudeste.
Considerando -se a avaliação da dificuldade das questões do Componente de Conhecimento Específico da prova, de acordo com o desempenho dos estudantes, não se observa diferença estatisticamente significativa.
Em todos os quartos, a proporção dos que classificaram a parte específica como difícil ou muito difícil variou de 26,7 \% ( 4º quarto ) a 31,3 \% ( 2º quarto ), numa tendência decrescente a partir do 2º quarto.
A alternativa modal para a Questão 2 foi o médio, com 57,0 \% no 1º quarto e 66,3 \% tanto no 3º quanto no 4º quarto. 4.2 EXTENSÃO DA PROVA EM RELAÇÃO AO TEMPO TOTAL Indagados quanto à extensão da prova, em relação ao tempo total oferecido para a sua resolução ( Questão 3 ), os estudantes apontaram, com maior incidência, a alternativa que considerava a extensão adequada, para todas as agregações consideradas ( Gráfico 4.5, Gráfico 4.6, e, no Anexo II, a Tabela II.3 ).
O percentual de alunos que responderam ser a extensão da prova adequada foi de 66,1 \%.
Já 29,9 \% dos inscritos e presentes consideraram que o exame foi longo ou muito longo e 4,1 \% o avaliaram como curto ou muito curto.
Entre as Grandes Regiões a proporção daqueles que avaliaram a prova como longa ou muito longa em relação ao tempo total destinado à sua resolução variou bastante : de 9,1 \% na região Norte a 34,0 \% na região Nordeste.
Observa -se no Gráfico 4.6 que a diferença entre as regiões Norte e Nordeste é estatisticamente significativa.
Considerando -se o desempenho dos alunos, nota -se que independente do quarto de desempenho entre 63 \% e 70 \% dos estudantes considerou a extensão da prova adequada.
E o comportamento em relação às demais opções desta pergunta não difere muito entre os quartos de desempenho.
No Gráfico 4.6, pode -se constatar que para os quatro quartos de desempenho a proporção de estudantes que consideraram a prova longa ou muito longa, em relação ao tempo total destinado à sua resolução, varou de 27,5 \% no quarto superior a 31,3 \% no 2º quarto, sem diferenças estatisticamente significativas.
A partir do 2º quarto nota -se uma tendência decrescente em relação ao aumento de desempenho. 4.3 COMPREENSÃO DOS ENUNCIADOS DAS QUESTÕES 4.3.1 Componente de Formação Geral Com relação aos enunciados das questões do Componente de Formação Geral ( Questão 4 ), as opiniões foram positivas, já que 76,6 \% dos alunos avaliados consideraram os enunciados de todas ou da maioria das questões claros e objetivos ( Gráfico 4.7, Gráfico 4.8, e, no Anexo II, a Tabela II.4 ).
Na análise regional, a percentagem de estudantes que avaliaram que todos ou a maioria dos enunciados das questões do Componente de Formação Geral \textbf{estavam} claros e objetivos variou de 71,4 \% na região Sudeste a 100,0 \% na região Centro-Oeste.
A análise das percepções dos estudantes sobre a clareza e objetividade dos enunciados permite afirmar que todos ou a maioria dos enunciados de questões relativas ao Componente de Formação Geral foram considerados claros e objetivos para a maior parte dos respondentes ( maior do que 71 \% em todas as regiões e maior do que 73 \% para todos os quartos de desempenho ).
Segundo o desempenho, observa -se que há uma tendência de crescimento da proporção dos que afirmaram que todos ou a maioria dos enunciados de questões eram claros e objetivos, sem diferenças estatisticamente significativas.
No quarto de desempenho inferior a proporção é de 73,3 \% e no de desempenho superior chega a 79,4 \%, sendo um pouco maior no 3º quarto, 79,8 \%. 4.3.2 Componente de Conhecimento Específico Com relação aos enunciados das questões do Componente de Conhecimento Específico da prova, para 80,9 \% dos estudantes avaliados da Área de Tecnologia em Saneamento Ambiental a clareza e a objetividade ( Questão 5 ) \textbf{estavam} presentes em todas ou na maioria das questões ( Gráfico 4.9, Gráfico 4.10, e no Anexo II, a Tabela II.5 ).
A maioria dos estudantes de todas as Grandes Regiões brasileiras considerou claros e objetivos todos ou a maioria dos enunciados das questões do Componente de Conhecimento Específico da prova, percentual sempre maior do que 76 \%.
As diferenças entre as duas menores proporções ( 76,6 \% da região Nordeste e 79,8 \% da região Sudeste ) e as duas maiores ( 97,0 \% da região centro-Oeste e 100,0 \% da região Sul ) são estatisticamente significativas.
A proporção de estudantes que consideraram os enunciados das questões claros e objetivos apresenta uma tendência crescente em relação ao aumento de desempenho : mais elevada no quarto superior ( 85,3 \% ) se comparada ao quarto inferior de desempenho ( 77,2 \% ).
No entanto, esta diferença não é estatisticamente significativa. 4.4 SUFICIÊNCIA DAS INFORMAÇÕES/INSTRUÇÕES FORNECIDAS Ao avaliarem as informações/instruções fornecidas para a resolução das questões ( Questão 6 ), 86,8 \% dos respondentes da Área de Tecnologia em Saneamento Ambiental de todo o Brasil afirmaram que estas eram até excessivas ou suficientes em todas ou na maioria das questões ( Gráfico 4.11, Gráfico 4.12, e, no Anexo II, a Tabela II.6 ).
A distribuição de respondentes que consideraram as informações/instruções fornecidas até excessivas ou suficientes em todas ou na maioria das questões pelas Grandes Regiões foi sempre superior a 80 \%, chegando a 100,0 \% na região Sul.
Levando -se em conta o desempenho dos participantes, as diferenças entre as proporções dos que avaliaram as informações/instruções como até excessivas ou suficientes em todas ou na maioria das questões, não são estatisticamente significativas, como mostra o Gráfico 4.12.
O percentual de participantes que emitiram tal opinião foi mais elevado no quarto superior ( 92,2 \% ), percentual superior à média nacional ( 86,8 \% ).
Já no quarto inferior, a suficiência das informações/instruções declarada como até excessiva em todas ou na maioria das questões foi percebida por 82,0 \% dos respondentes.
Observa -se uma tendência crescente conforme aumenta o desempenho. 4.5 DIFICULDADE ENCONTRADA AO RESPONDER À PROVA \textbf{Perguntados} sobre as dificuldades com as quais se depararam ao responder à prova ( Questão 7 ), 12,7 \% dos estudantes apontaram o desconhecimento do conteúdo.
Para 42,1 \%, a forma diferente de abordagem do conteúdo foi indicada como dificuldade.
Já a falta de motivação para fazer a prova foi apontada por 27,4 \% dos respondentes.
Considerando -se todo o Brasil, 14,2 \% dos respondentes afirmaram que não tiveram qualquer tipo de dificuldade para responder à prova ( Tabela II.7 no Anexo II ).
Os Gráficos 4.13 e 4.14 apresentam os percentuais de estudantes que apontaram o desconhecimento do conteúdo como dificuldade percebida ao responder à prova.
Na análise por Grandes Regiões, o maior percentual de inscritos e presentes que apontaram o desconhecimento do conteúdo como dificuldade ao responder à prova ocorreu na região Norte, 21,2 \% dos inscritos e presentes.
Já na região Centro-Oeste nenhum aluno escolheu esta opção de resposta.
A forma diferente de abordagem do conteúdo foi a escolha modal dos estudantes, com percentuais que variaram de 27,3 \% ( região Norte ) a 55,3 \% ( Nordeste ).
O percentual de alunos que citou a falta de motivação como dificuldade variou de 18,4 \% ( região Nordeste ) a 35,0 \% ( região Sudeste ).
Os que declararam não ter qualquer dificuldade para responder à prova variaram de 7,1 \% na região Nordeste a 33,3 \% na Sul.
Com relação aos quartos de desempenho, o desconhecimento do conteúdo, foi a opção escolhida por 13,0 \% dos estudantes do quarto inferior aumentando para 17,0 \% no 2º quarto.
A partir do 2º quarto há uma tendência decrescente da proporção de alunos com tal opinião, que chega a 9,0 \% dos estudantes do quarto superior.
A alternativa modal para os alunos, quando agregados pelos quartos de desempenho, foi a forma diferente de abordagem do conteúdo : 47,0 \% nos quartos extremos, e 38,4 \% e 36,5 \%, respectivamente no 2º e 3º quarto. 4.6 CONTEÚDOS DAS QUESTÕES OBJETIVAS DA PROVA Ao analisarem os conteúdos das questões objetivas da prova ( Questão 8 ), um percentual muito pequeno dos estudantes avaliados, apenas 3,8 \%, afirmou que não estudou ainda a maioria desses conteúdos ( Gráficos 4.15, Gráfico 4.16, e a Tabelas II.8 no Anexo II ).
A maioria ( 71,3 \% ) afirmou ter estudado e aprendido muitos ou todos os conteúdos avaliados.
Na análise por Grande Região, a proporção de respondentes que escolheram a opção não estudou ainda a maioria desses conteúdos, foi pequena.
Observa -se que nas regiões Sudeste ( 4,0 \% ) e Nordeste ( 4,3 \% ), apesar de pequenas, as proporções são maiores que a média nacional ( 3,8 \% ).
Não se observa diferença estatisticamente significativa entre as regiões.
Em todas as regiões, a maioria dos presentes afirmou ter estudado e aprendido muitos ou todos os conteúdos, com proporções variando entre 63,6 \% na região Norte e 78,8 \% na Centro-Oeste.
Considerando -se separadamente as opiniões de estudantes dos quatro quartos de desempenho, observa -se que, no quarto inferior, 7,0 \% ofereceram como resposta que não estudou ainda a maioria desses conteúdos, 2,0 \% os do quarto superior com a mesma resposta.
No entanto esta diferença não é estatisticamente significativa.
Tendo em conta o quarto superior, 87,3 \% dos alunos afirmaram ter estudado e aprendido muitos ou todos os conteúdos. 4.7 TEMPO GASTO PARA CONCLUIR A PROVA Ao responderem sobre o tempo de conclusão da prova ( Questão 9 ), mais da metade dos estudantes ( 72,3 \% ) afirmou ter gasto entre duas e quatro horas ( Gráfico 4.17, Gráfico 4.18 e, no Anexo II, a Tabela II.9 ).
Considerando -se as cinco Grandes Regiões brasileiras, o percentual dos que utilizaram entre duas e quatro horas para finalizar a prova ficou abaixo do nacional nas regiões Nordeste ( 69,3 \% ) e Sudeste ( 70,6 \% ).
Na região Norte a proporção dos que dispensaram entre duas e quatro horas para concluir a prova atingiu 87,9 \% dos alunos.
Uma vez sendo analisadas as alternativas escolhidas pelos estudantes dos diferentes quartos de desempenho, observa -se que uma maior proporção de participantes do 3º quarto declarou ter gasto entre duas e quatro horas para concluir a prova.
Nos dois quartos inferiores tal proporção foi de 78,8 \% e 52,0 \%, respectivamente, no 1º e 2º quartos.
As diferenças entre o primeiro e o segundo quartos são estatisticamente significativas.
As diferenças entre os quartos evidenciam uma tendência crescente do 1º ao 3º, com diferenças estatisticamente significativas entre o resultado do quarto inferior em relação aos demais.

% matched lemmas: cada, estar, menor, perguntar
\end{textsample}
